\begin{abstract}
\noindent
Decentralized-finance (DeFi) liquidations record forced deleveraging
transaction by transaction in public data (a granularity unavailable in
centralized venues, where such deleveraging is essentially unobservable).
The leverage-cycle literature predicts that such forced selling should
amplify downside tail risk. We test this prediction for Ethereum (ETH)
using quantile local projections on \Nobs{} hourly observations (2021--2025).
A naive specification appears to confirm it: the response of the first return
percentile to a liquidation shock is roughly \RatioTailMedImpact$\times$ the
median response on impact, deepens with the horizon, and survives
conventional robustness checks with Bonferroni correction in full. A battery of
diagnostics overturns this reading. A symmetric sign-flip placebo reproduces
the entire left--right tail gap at every horizon; the response is equally
strong in the \emph{upper} tail; the same signature appears when Bitcoin is
substituted as the outcome (at $\BtcOverEthHtwelve$--$\BtcOverEthHzero\times$ the magnitude); and dual
permutation nulls show the long-horizon coefficient is genuine rather than a
mechanical artifact of overlapping cumulative returns. Net of volatility, an
equivalence test \emph{bounds} the downside-specific tail-violation asymmetry, at the one-percent tail, below economically negligible thresholds under all but the strictest calibration of the single detectable residual. What survives is a
\emph{symmetric} volatility channel: liquidations carry out-of-sample
predictive content for both moderate return tails at short horizons, incremental to a
nested quantile benchmark though not a substitute for a dedicated volatility
model.
The exercise yields a portable methodological caution. In quantile local
projections on overlapping cumulative returns, symmetric heteroskedasticity
masquerades as downside fragility.
\end{abstract}